\documentclass[../../../include/open-logic-section]{subfiles}

\begin{document}

\olfileid{sth}{replacement}{absinf}
\olsection{تعويض او «مطلق نامتناهيت»}

اوس \limofsize{} شاته پرېږدو او د تعويض د توجيې د نورو
(دننيو) هڅو يوه ډله څېړو۔ دغه هڅې دا خبره جدي نيسي چې سټونه په
پړاوونو کښې جوړېږي۔

کله چې مو لومړے تکراري بهير بيان کړ، داسې اصول مو وړاندې کړل چې په
هر پړاو کښې پېښېدونکے کار ئې روښانولو۔ دا \stageshier، \stagesord او
\stagesacc وو۔ وروسته مو داسې اصول ورزيات کړل چې د پړاوونو د شمېر
په اړه څه وايي: \stagessucc{} وايي چې د سټ جوړولو بهير هېڅکله نۀ
ختمېږي، او \stagesinf{} وايي چې بهير تر نامتناهيتم پړاو پورې رسېږي۔

دا مناسبه ده چې ووايو دغه وروستي دوه اصول ښايي له يوه لا پراخ اصل
څخه راووځي، لکه:
\begin{enumerate}
	\item[] \stagesinex. په مطلق ډول نامتناهي ډېر پړاوونه شته؛
	سلسله تر هغه حده لوړه ده چې امکان لري.
\end{enumerate}
ښکاره ده چې دا يو غير رسمي اصل دے۔ خو که دا د سټ له تجمعي-تکراري
تصور څخه سمدستي \emph{نۀ هم لازمېږي}، له هغه سره خو \emph{سمون}
لري۔ لږ تر لږه، او د \limofsize پر خلاف، دا تصور ساتي چې سټونه
يو پړاو په بل پسې جوړېږي۔

اوس هيله دا ده چې له \stagesinex{} څخه د تعويض توجيه ترلاسه کړو۔
راځئ وګورو چې دا څنګه کېدے شي۔

په \olref[ordinals][idea]{sec} کښې مو وليدل چې داسې ښه ترتيب جوړول
اسانه دي چې د مفهوم له مخې بايد له $\omega+\omega$ سره آيزومورف
وي۔ په بل عبارت، په اسانه داسې پړاو په پړاو تکراري بهير تصورولے شو
چې د ترتيب ډول ئې د مفهوم له مخې $\omega+\omega$ وي۔ نو که مو
\stagesinex منلے وي، خامخا بايد ومنو چې د سلسلې لږ تر لږه
$\omega+\omega$-م پړاو، يعنې $V_{\omega+\omega}$، شته؛ ځکه سلسله
تر دې ځايه \emph{غځېدې} شي۔

دا فکر داسې عامېږي: د هر ښه ترتيب دپاره بايد د تکراري سلسلې د
جوړولو بهير لږ تر لږه د هماغه ښه ترتيب تر اوږدوالي ورسېږي۔ دا خبره
په دې هم تضمينولے شو چې \olref[ordinals][ordtype]{thmOrdinalRepresentation}
د يو \emph{اصل} په توګه واخلو۔ نو هر ښه ترتيب به د فون نويمن له
کوم ترتيبي عدد سره آيزومورف وي۔ څرنګه چې هر فون نويمن ترتيبي عدد د
خپلې رتبې سره برابر وي، \olref[ordinals][ordtype]{thmOrdinalRepresentation}
به راته ووايي چې کله هم په خپله سټ تيورۍ کښې يو ښه ترتيب بيانولے
شو، د سټ جوړولو تکراري بهير بايد تر هغه ښه ترتيب هاخوا ولاړ شي۔

دا فکر بېشکه د \stagesinex يوه پايله ښکاري۔ له بده مرغه، که مو موخه
له دې فکره د تعويض راايستل وي، يوه ساده تخنيکي خنډ مخې ته راځي:
تعويض له \olref[ordinals][ordtype]{thmOrdinalRepresentation} څخه په
کلکه قوي دے۔ (دا خبره \citet[\S13.2]{Potter2004} کړې ده؛ موږ به ئې
په \olref[sth][replacement][finiteaxiomatizability]{sec} کښې ثابته کړو۔)

نو که \stagesinex{} داسې وپېژنو چې تعويض ترې راووځي، دا اصل
\emph{يوازې} دومره نۀ شي ويلے چې سلسله تر هر ښه ترتيب هاخوا غځېږي۔
بايد تر دې لا قوي خبره وکړي۔ د همدې دپاره \cite{Shoenfield:AST} د
فکر يوه ډېره طبيعي پياوړې بڼه وړاندې کړه: سلسله له هېڅ سټ سره
\emph{هم‌پايه} نۀ ده۔\footnote{ښايي G\"odel ورته فکر وړاندې کړے وي؛
\citet[p.~223]{Potter2004} وګورئ۔ د G\"odel او
\citeauthor{Shoenfield:AST} د بحث دپاره \citet[90--5]{Incurvati2020}
وګورئ۔} په لږ تفصيل: که $\tau$ داسې نگاشت وي چې سټونه د سلسلې
پړاوونو ته لېږي، د هر سټ $A$ انځور د $\tau$ لاندې ټوله سلسله نۀ شي
پوښلے۔ په بل عبارت، په شېمايي ډول:
\begin{enumerate}
	\item[] \stagescofin. که $A$ يو سټ وي او $\tau(x)$ د هر
	$x \in A$ دپاره يو پړاو وي، نو داسې پړاو شته چې د هر
	$\tau(x)$ نه وروسته راځي، د هر $x \in A$ دپاره.
\end{enumerate}
ښکاره ده چې $\ZF$ د \stagescofin يوه مناسبه رسمي بڼه ثابتوي۔
برعکس، په غير رسمي استدلال سره ښودلای شو چې \stagescofin{} تعويض
توجيه کوي۔\footnote{د \stagescofin د کومې رسمي بڼې نه تعويض ثابتول
به ګران وي، ځکه $\Z$ په خپله د پړاوونو د تعريف دپاره کافي قوي نۀ
ده؛ نو څرګنده نۀ ده چې \stagescofin به څنګه رسمي کړو۔ يوه لاره دا
ده چې د $\LT$ په کوم پراخ نظام کښې کار وکړو، لکه په
\olref[sth][replacement][extrinsic]{sec} کښې چې بحث شوے دے۔}
فرض کړئ $(\forall x \in A)\lexists![y][\phi(x,y)]$۔ نو د هر
$x \in A$ دپاره $\sigma(x)$ هغه $y$ وټاکئ چې $\phi(x,y)$ پوره
کوي، او $\tau(x)$ هغه پړاو وټاکئ چې $\sigma(x)$ په لومړي ځل پکښې
جوړېږي۔ د \stagescofin له مخې داسې پړاو $V$ شته چې
$(\forall x \in A)\tau(x)\in V$۔ اوس، څرنګه چې هر
$\tau(x) \in V$ دے او $\sigma(x) \subseteq \tau(x)$، د بېلونې
په اصل سره
$\Setabs{y
\in V}{(\exists x \in A)\sigma(x) = y} = \Setabs{y}{(\exists x \in
A)\phi(x,y)}$ ترلاسه کولے شو۔

\begin{prob}
	\stagescofin{} د $\ZF$ دننه په رسمي ډول بيان کړئ.
\end{prob}

نو \stagescofin{} تعويض توجيه کوي۔ او لږ تر لږه دا معقوله ښکاري
چې \stagesinex{}، \stagescofin توجيه کوي۔ فرض کړئ \stagescofin{}
نۀ پوره کېږي۔ نو سلسله له کوم سټ~$A$ سره هم‌پايه ده: داسې نگاشت
$\tau$ لرو چې د هر پړاو~$S$ دپاره کوم $x \in A$ شته چې
$S \in \tau(x)$۔ په دې حالت کښې د سلسلې د ادعا شوي «مطلق نامتناهيت»
د اندازه کولو يوه لاره لرو: ټوله سلسله د $\tau$ د $A$ پر غړو د
تطبيق د قيمتونو سټ په وسيله \emph{پوښل کېږي}۔ دا د سلسلې د «مطلق نامتناهيت» فکر کمزورے
کوي۔ په معکوس استدلال: \stagesinex{}، \stagescofin لازموي، او هغه
بيا تعويض توجيه کوي۔

دا د تعويض د دننۍ توجيې يوه رښتيا هيله‌بښونکې هڅه ده۔ خو په پای کښې
دا کار کوي که نۀ، پرېکړه ئې تاسو ته پرېږدو۔

\end{document}
